\hypertarget{ux57faux4e8ecux4e0eqt6ux7684ux4ea4ux4e92ux5f0fux6587ux5b57ux5192ux9669ux6e38ux620fux8bbeux8ba1ux4e0eux5b9eux73b0-ux94a2ux94c1ux610fux5fd7ux7b2cux4e09ux5e1dux56fdux7684ux9ec4ux660f}{%
\section{基于C++与Qt6的交互式文字冒险游戏设计与实现 ------
《钢铁意志：第三帝国的黄昏》}\label{ux57faux4e8ecux4e0eqt6ux7684ux4ea4ux4e92ux5f0fux6587ux5b57ux5192ux9669ux6e38ux620fux8bbeux8ba1ux4e0eux5b9eux73b0-ux94a2ux94c1ux610fux5fd7ux7b2cux4e09ux5e1dux56fdux7684ux9ec4ux660f}}

\begin{center}\rule{0.5\linewidth}{0.5pt}\end{center}

\hypertarget{ux6458ux8981}{%
\subsection{摘要}\label{ux6458ux8981}}

本文介绍了一款基于C++17与Qt6框架开发的交互式文字冒险游戏------《钢铁意志：第三帝国的黄昏》。游戏以第二次世界大战欧洲战场为背景，采用MVC分层架构，实现了包含六大职业系统、五大战役场景、D100骰子判定系统、叙事标志分支系统、打字机效果以及JSON持久化存档等核心功能。游戏共编写约3000行C++代码，涵盖28个源文件，构建了一个从1940年黄色方案到1945年柏林战役的完整叙事体系。本文从项目背景、技术架构、核心系统实现、界面设计以及构建测试等方面，系统阐述了该游戏的完整开发过程。

\textbf{关键词：} C++17；Qt6；文字冒险游戏；MVC架构；二战叙事

\begin{center}\rule{0.5\linewidth}{0.5pt}\end{center}

\hypertarget{ux7eeaux8bba}{%
\subsection{1. 绪论}\label{ux7eeaux8bba}}

\hypertarget{ux5b9eux4e60ux80ccux666f}{%
\subsubsection{1.1 实习背景}\label{ux5b9eux4e60ux80ccux666f}}

本次实习是面向对象程序设计实践课程的组成部分。课程要求运用C++语言及面向对象思想，完成一个具有一定规模的实际项目。游戏项目因其天然的交互性、状态复杂性和多媒体需求，成为检验面向对象程序设计能力的理想载体。

\hypertarget{ux9009ux9898ux52a8ux673a}{%
\subsubsection{1.2 选题动机}\label{ux9009ux9898ux52a8ux673a}}

战争题材的游戏作品繁多，但大多聚焦于战术操作或战略模拟，鲜有作品关注战争本身的残酷与荒诞。本项目的设计初衷是创作一部``文字悲剧''------让玩家通过德军士兵的视角，亲历战争中的道德困境、战友消亡与帝国崩溃，从而反思战争的本质。技术层面上，选题涵盖了Qt信号槽通信、JSON数据序列化、状态机管理、随机数判定系统等多个面向对象编程核心概念，具有良好的教学实践价值。

\hypertarget{ux9879ux76eeux76eeux6807}{%
\subsubsection{1.3 项目目标}\label{ux9879ux76eeux76eeux6807}}

本项目旨在实现一个具备以下特性的交互式文字冒险游戏：

\begin{enumerate}
\def\labelenumi{\arabic{enumi}.}
\tightlist
\item
  \textbf{完整的叙事体系}：涵盖五个历史战役，每个战役包含10-28个分支故事节点
\item
  \textbf{多元的游戏系统}：职业选择、骰子战斗、道德标志、存档回溯
\item
  \textbf{专业的用户界面}：军工战损风主题，打字机逐字渲染效果
\item
  \textbf{良好的软件架构}：分层解耦、信号驱动、可扩展场景系统
\end{enumerate}

\begin{center}\rule{0.5\linewidth}{0.5pt}\end{center}

\hypertarget{ux9879ux76eeux6982ux8ff0}{%
\subsection{2. 项目概述}\label{ux9879ux76eeux6982ux8ff0}}

\hypertarget{ux6e38ux620fux80ccux666f}{%
\subsubsection{2.1 游戏背景}\label{ux6e38ux620fux80ccux666f}}

《钢铁意志：第三帝国的黄昏》以第二次世界大战欧洲战场为叙事框架。玩家从1940年德国发动黄色方案（Fall
Gelb）开始，依次经历不列颠空战、大西洋海战、斯大林格勒战役，最终在1945年柏林战役中迎来第三帝国的覆灭。游戏的历史改编基于《第三帝国的兴亡》《斯大林格勒》等战史著作，以及《我们的父辈》《帝国的毁灭》等影视作品。

\hypertarget{ux6838ux5fc3ux7279ux8272}{%
\subsubsection{2.2 核心特色}\label{ux6838ux5fc3ux7279ux8272}}

\begin{itemize}
\tightlist
\item
  \textbf{六大职业路线}：步兵、坦克兵、战斗机飞行员、轰炸机飞行员、潜艇驾驶员、战列舰水兵，不同职业对应不同的初始战役与骰子判定加成
\item
  \textbf{五大战役}：每个战役基于真实历史事件改编，包含多条分支路径与总计15种以上结局
\item
  \textbf{D100骰子战斗系统}：随机判定结合职业加成，决定战斗选项的成败
\item
  \textbf{叙事标志系统}：玩家的道德选择会留下永久标志（如``spared\_village''``saved\_survivors''），影响后续选项可见性
\item
  \textbf{打字机效果}：逐字打印的故事叙述，增强氛围沉浸感
\item
  \textbf{存档系统}：1个自动存档位加3个手动存档位，JSON格式持久化存储
\end{itemize}

\begin{center}\rule{0.5\linewidth}{0.5pt}\end{center}

\hypertarget{ux6280ux672fux67b6ux6784}{%
\subsection{3. 技术架构}\label{ux6280ux672fux67b6ux6784}}

\hypertarget{ux6280ux672fux9009ux578b}{%
\subsubsection{3.1 技术选型}\label{ux6280ux672fux9009ux578b}}

\begin{longtable}[]{@{}lll@{}}
\toprule
层次 & 技术 & 选型理由 \\
\midrule
\endhead
语言 & C++17 &
课程要求，性能优越，支持现代C++特性如\texttt{std::clamp}、\texttt{std::unique\_ptr} \\
GUI框架 & Qt 6 (Core/Widgets/Multimedia) & 成熟的跨平台C++
GUI框架，内置信号槽机制天然适合事件驱动架构 \\
构建系统 & CMake 3.16+ &
跨平台构建，与Qt6深度集成，支持AUTOMOC/AUTORCC \\
音频 & QMediaPlayer + QAudioOutput & Qt6
Multimedia模块原生支持，实现背景音乐播放与循环 \\
数据序列化 & QJsonDocument/QJsonObject &
Qt内置JSON支持，无需第三方依赖 \\
样式 & Qt Style Sheets (QSS) &
CSS语法，支持选择器、伪类和动态属性，实现暗色军工主题 \\
\bottomrule
\end{longtable}

\hypertarget{ux9879ux76eeux7ed3ux6784}{%
\subsubsection{3.2 项目结构}\label{ux9879ux76eeux7ed3ux6784}}

项目严格遵循MVC（Model-View-Controller）分层架构，通过Qt信号槽机制实现层间解耦通信：

\begin{verbatim}
Hearts of iron/
├── main.cpp                    # 应用程序入口，加载全局样式
├── CMakeLists.txt              # CMake构建配置
├── core/                       # 核心逻辑层（Model + Controller）
│   ├── GameState.h             # 全局枚举（职业/场景/节点类型）与辅助函数
│   ├── StoryNode.h             # 故事节点（StoryNode）与选项（Choice）数据结构
│   ├── Player.h/.cpp           # 玩家数据模型，含状态管理与JSON序列化
│   ├── GameEngine.h/.cpp       # 核心游戏引擎，中央调度器与信号发射
│   ├── ScenarioBase.h/.cpp     # 场景抽象基类，定义节点注册与查找接口
│   ├── SaveManager.h/.cpp      # 存档管理器，JSON文件读写
│   └── MusicPlayer.h/.cpp      # 背景音乐管理，键-路径映射
├── scenarios/                  # 五大战役场景（Model数据定义）
│   ├── FallGelbScenario.h/.cpp     # 黄色方案 (1940) — 28个节点
│   ├── BritainScenario.h/.cpp      # 不列颠空战 (1940) — 27个节点
│   ├── WolfPackScenario.h/.cpp     # 群狼海战 (1941-1943) — 25个节点
│   ├── StalingradScenario.h/.cpp   # 斯大林格勒 (1942-1943) — 27个节点
│   └── BerlinScenario.h/.cpp       # 柏林战役 (1945) — 22个节点
├── ui/                         # 用户界面层（View）
│   ├── MainWindow.h/.cpp           # 主窗口，QStackedWidget页面管理与信号路由
│   ├── MenuWidget.h/.cpp           # 主菜单界面
│   ├── CharacterCreateWidget.h/.cpp # 角色创建界面
│   ├── GameWidget.h/.cpp           # 游戏主界面，含打字机效果实现
│   └── SaveLoadDialog.h/.cpp       # 存档/读档模态对话框
└── resources/                  # 资源文件
    ├── resources.qrc               # Qt资源索引文件
    └── style.qss                   # 全局QSS样式表（军工战损风）
\end{verbatim}

\hypertarget{ux67b6ux6784ux8bbeux8ba1}{%
\subsubsection{3.3 架构设计}\label{ux67b6ux6784ux8bbeux8ba1}}

项目的核心设计思想是``引擎集中调度，信号异步通知，UI被动响应''。整个数据流是单向的：玩家点击按钮
→ UI发射信号 → MainWindow转发 →
GameEngine处理（骰子判定/标志设置/HP计算） → 引擎发射信号 →
MainWindow接收 →
刷新UI各项组件。这种设计使得核心逻辑层完全不依赖UI层，可以独立进行单元测试，也便于未来更换前端框架。

各层职责如下：

\begin{itemize}
\tightlist
\item
  \textbf{数据模型层（StoryNode/Player/ScenarioBase）}：定义故事节点图、玩家状态、场景节点集合等核心数据结构
\item
  \textbf{逻辑控制层（GameEngine/SaveManager/MusicPlayer）}：处理选择流程、骰子判定、状态变更、存档读写和音乐播放
\item
  \textbf{用户界面层（MenuWidget/CharacterCreateWidget/GameWidget/SaveLoadDialog）}：负责布局渲染、打字机动画、用户交互和视觉反馈
\item
  \textbf{路由层（MainWindow）}：使用QStackedWidget管理页面切换，连接并转发各层信号，维护应用程序生命周期
\end{itemize}

\begin{center}\rule{0.5\linewidth}{0.5pt}\end{center}

\hypertarget{ux6838ux5fc3ux7cfbux7edfux5b9eux73b0}{%
\subsection{4.
核心系统实现}\label{ux6838ux5fc3ux7cfbux7edfux5b9eux73b0}}

\hypertarget{ux573aux666fux56feux4e0eux53d9ux4e8bux7cfbux7edf}{%
\subsubsection{4.1
场景图与叙事系统}\label{ux573aux666fux56feux4e0eux53d9ux4e8bux7cfbux7edf}}

场景系统的核心是场景图（Scene
Graph）------一个有向图结构，其中节点（\texttt{StoryNode}）代表叙事片段，边（\texttt{Choice::nextNodeId}）代表玩家选择导致的叙事跳转。所有场景继承自抽象基类
\texttt{ScenarioBase}，必须实现三个纯虚方法：

\begin{Shaded}
\begin{Highlighting}[]
\KeywordTok{virtual} \DataTypeTok{void}\NormalTok{ initialize}\OperatorTok{()} \OperatorTok{=} \DecValTok{0}\OperatorTok{;}                              \CommentTok{// 构建节点图}
\KeywordTok{virtual}\NormalTok{ ScenarioId scenarioId}\OperatorTok{()} \AttributeTok{const} \OperatorTok{=} \DecValTok{0}\OperatorTok{;}                  \CommentTok{// 返回场景枚举}
\KeywordTok{virtual} \ExtensionTok{QString}\NormalTok{ startNodeId}\OperatorTok{(}\NormalTok{PlayerClass cls}\OperatorTok{)} \AttributeTok{const} \OperatorTok{=} \DecValTok{0}\OperatorTok{;}    \CommentTok{// 职业起点}
\end{Highlighting}
\end{Shaded}

每个场景在构造函数中调用 \texttt{initialize()}，通过多次调用
\texttt{addNode(StoryNode)} 构建内部节点图。节点以
\texttt{QMap\textless{}QString,\ StoryNode\textgreater{}}
形式存储，支持O(log n)查找。故事节点包含三种类型：

\begin{itemize}
\tightlist
\item
  \textbf{Narrative（叙事节点）}：纯文字叙述，仅提供``继续''按钮，直接跳转到
  \texttt{nextNodeId}
\item
  \textbf{Choice（选择节点）}：包含多个 \texttt{Choice}
  选项，每个选项可设置战斗判定、前置标志条件、职业限制等
\item
  \textbf{Ending（终止节点）}：标记 \texttt{isVictory} 或
  \texttt{isDefeat}，触发战役结束逻辑
\end{itemize}

选项结构（\texttt{Choice}）是系统的核心，包含以下关键字段：

\begin{itemize}
\tightlist
\item
  \texttt{isCombat} 与
  \texttt{combatThreshold}：标记战斗选项及其骰子阈值
\item
  \texttt{bonusClasses}：享受+20加成的职业列表
\item
  \texttt{requiredFlags} /
  \texttt{grantedFlags}：前置标志条件与选择后授予的标志
\item
  \texttt{classRestricted} / \texttt{allowedClasses}：职业限制
\item
  \texttt{successNodeId} /
  \texttt{failureNodeId}：成功与失败的不同跳转路径
\end{itemize}

这种设计使得每个选项都可以承载丰富的分支逻辑，同时保持数据结构的简洁性。

\hypertarget{ux9ab0ux5b50ux6218ux6597ux7cfbux7edf}{%
\subsubsection{4.2
骰子战斗系统}\label{ux9ab0ux5b50ux6218ux6597ux7cfbux7edf}}

战斗判定采用D100随机数系统，是游戏最核心的交互机制。实现代码如下：

\begin{Shaded}
\begin{Highlighting}[]
\DataTypeTok{int}\NormalTok{ GameEngine}\OperatorTok{::}\NormalTok{rollDice}\OperatorTok{(}\AttributeTok{const} \ExtensionTok{QList}\OperatorTok{\textless{}}\NormalTok{PlayerClass}\OperatorTok{\textgreater{}} \OperatorTok{\&}\NormalTok{bonusClasses}\OperatorTok{)} \AttributeTok{const} \OperatorTok{\{}
    \DataTypeTok{int}\NormalTok{ roll }\OperatorTok{=} \ExtensionTok{QRandomGenerator::}\NormalTok{global}\OperatorTok{(){-}\textgreater{}}\NormalTok{bounded}\OperatorTok{(}\DecValTok{1}\OperatorTok{,} \DecValTok{101}\OperatorTok{);} \CommentTok{// [1, 100]}
    \ControlFlowTok{if} \OperatorTok{(}\NormalTok{bonusClasses}\OperatorTok{.}\NormalTok{contains}\OperatorTok{(}\VariableTok{m\_player}\OperatorTok{.}\NormalTok{playerClass}\OperatorTok{))}
\NormalTok{        roll }\OperatorTok{+=} \DecValTok{20}\OperatorTok{;}                       \CommentTok{// 职业匹配 → +20 加成}
    \ControlFlowTok{return} \FunctionTok{qBound}\OperatorTok{(}\DecValTok{1}\OperatorTok{,}\NormalTok{ roll}\OperatorTok{,} \DecValTok{120}\OperatorTok{);}          \CommentTok{// 钳制范围 [1, 120]}
\OperatorTok{\}}
\end{Highlighting}
\end{Shaded}

判定逻辑简洁但富有层次：基础随机值1-100，若玩家职业在选项的
\texttt{bonusClasses} 列表中则获得+20加成，最终值钳制在 {[}1, 120{]}
范围内。判定成功条件为最终投掷值 ≥ 选项的 \texttt{combatThreshold}。

阈值设计体现了游戏难度曲线：标准战斗阈值50（无加成成功率51\%，有加成71\%），高强度战斗阈值55-65，极限行动阈值70，最终死战阈值90。这种梯度设计使得游戏早期的选择相对宽容，而后期（如斯大林格勒的最后抵抗）则让玩家切实感受到生存的艰难。

战斗成功时执行 \texttt{applyChoice()}
应用HP/士气变化与标志授予，跳转到成功节点；失败时应用
\texttt{failHpDelta} 与 \texttt{failMoraleDelta}，若HP或士气降至零则触发
\texttt{isDead()}
判定，导航到失败结局节点。这一完整的战斗闭环包含了随机性、策略性（玩家选择影响阈值）和后果反馈。

\hypertarget{ux53d9ux4e8bux6807ux5fd7ux7cfbux7edf}{%
\subsubsection{4.3
叙事标志系统}\label{ux53d9ux4e8bux6807ux5fd7ux7cfbux7edf}}

叙事标志系统是实现故事分支的核心机制。玩家在关键道德选择中做出的决定会以
\texttt{QSet\textless{}QString\textgreater{}}
形式永久存储。标志的典型应用场景是``前置条件检查''------某些选项通过
\texttt{requiredFlags} 字段声明只有在玩家拥有特定标志时才可见。例如：

\begin{itemize}
\tightlist
\item
  在黄色方案中选择了``禁止骚扰平民''获得 \texttt{spared\_village}
  标志后，后续场景中的人道主义选项才会出现
\item
  在不列颠空战中拒绝轰炸平民目标获得 \texttt{refused\_bomb}
  标志，影响后续的军官对话文本
\end{itemize}

UI层在执行选项过滤时会遍历当前节点的所有选项，跳过职业不匹配或标志不满足的选项，不生成对应按钮。这种设计使得分支逻辑完全由数据驱动，场景设计者只需在节点数据中声明条件，无需修改任何控制代码。

\hypertarget{ux5b58ux6863ux7cfbux7edf}{%
\subsubsection{4.4 存档系统}\label{ux5b58ux6863ux7cfbux7edf}}

存档系统采用JSON格式持久化玩家状态，通过 \texttt{SaveManager}
类管理。存储路径使用Qt标准机制
\texttt{QStandardPaths::writableLocation(QStandardPaths::AppDataLocation)}
确保跨平台兼容（Windows下为
\texttt{\%APPDATA\%/HeartsOfIronGame/saves/}），共1个自动存档位（Slot
0）和3个手动存档位（Slot 1-3）。

存档数据结构将玩家完整状态（姓名、职业、HP、士气、标志集合、当前场景与节点ID、已解锁场景列表）序列化为JSON对象，附加槽位编号、时间戳等元信息。读取时使用
\texttt{Player::fromJson()}
方法进行反序列化，所有字段均设有默认值容错（如 \texttt{hp}
默认100），确保数据文件部分损坏时仍可恢复。

值得注意的是，场景数据（节点文本、选项逻辑）不随存档序列化------这些是硬编码的常量图。这一设计决策简化了存档结构，也避免了场景数据变更导致的存档兼容性问题。自动存档在每次节点切换时触发（\texttt{onNodeChanged}中调用
\texttt{autoSave()}），确保玩家进度不丢失。

\hypertarget{ux6253ux5b57ux673aux6548ux679c}{%
\subsubsection{4.5 打字机效果}\label{ux6253ux5b57ux673aux6548ux679c}}

打字机效果是提升叙事沉浸感的关键UI特性，实现在 \texttt{GameWidget}
中。其核心机制如下：

\begin{itemize}
\tightlist
\item
  使用 \texttt{QTimer} 每20ms触发一次 \texttt{onTypeTimerTick()}，每次在
  \texttt{QTextBrowser} 光标末尾插入一个字符
\item
  打字期间禁用所有选项按钮和存档/读档按钮，防止玩家提前操作打断叙事节奏
\item
  通过在 \texttt{QTextBrowser} 及其 \texttt{viewport}
  上安装事件过滤器（\texttt{eventFilter}）拦截鼠标点击事件，实现``点击跳过''功能------点击后直接调用
  \texttt{displayFullText()} 渲染全部文字
\item
  打字过程中强制将滚动条滚动到底部，确保最新文字始终可见
\end{itemize}

这一实现细节体现了对用户体验的细致考量：既要营造沉浸感（不让玩家急于一扫而过），又要尊重玩家的自主权（允许跳过）。

\hypertarget{ux97f3ux4e50ux7cfbux7edf}{%
\subsubsection{4.6 音乐系统}\label{ux97f3ux4e50ux7cfbux7edf}}

音乐系统使用 \texttt{QMediaPlayer} + \texttt{QAudioOutput}
实现，采用键-路径映射管理音轨。每个故事节点通过 \texttt{musicKey}
字段声明需要播放的音乐。系统具有以下设计特点：

\begin{itemize}
\tightlist
\item
  \textbf{播放去重}：如果请求的键与当前键相同且正在播放，则跳过切换
\item
  \textbf{循环播放}：通过连接 \texttt{mediaStatusChanged} 信号监听
  \texttt{EndOfMedia} 状态并重置播放位置
\item
  \textbf{文件容错}：播放前检查文件是否存在（\texttt{QFileInfo::exists}），不存在则静默跳过，不影响核心逻辑
\item
  \textbf{静音支持}：\texttt{setMuted(true)}
  将音量设为零但不改变存储的音量值
\end{itemize}

\begin{center}\rule{0.5\linewidth}{0.5pt}\end{center}

\hypertarget{ux754cux9762ux8bbeux8ba1}{%
\subsection{5. 界面设计}\label{ux754cux9762ux8bbeux8ba1}}

\hypertarget{ux4e3bux9898ux98ceux683c}{%
\subsubsection{5.1 主题风格}\label{ux4e3bux9898ux98ceux683c}}

游戏全局采用``军工战损风''（Industrial
War-Worn）主题。设计语言的核心要素包括：

\begin{itemize}
\tightlist
\item
  \textbf{暗色层级}：背景 \texttt{\#121212} → 面板 \texttt{\#181818} →
  卡片 \texttt{\#1c1c1c} → 悬停 \texttt{\#242424}，四层灰色建立纵深
\item
  \textbf{强调色}：\texttt{\#aa3333} /
  \texttt{\#ff3333}（暗红，模拟血与铁锈），\texttt{\#ff9900}（橙色，关键信息标注）
\item
  \textbf{无圆角}：所有元素
  \texttt{border-radius:\ 0px}，体现军事工业的冷硬直线
\item
  \textbf{等宽回退字体}：\texttt{font-family} 末尾回退到
  \texttt{monospace}，营造电报/打字机感
\item
  \textbf{功能色语义化}：HP进度条使用血红（\texttt{\#992222}），士气进度条使用冷蓝（\texttt{\#226699}），战斗选项按钮使用红色调区分风险行动
\end{itemize}

所有视觉效果通过QSS样式表实现，无外部图片依赖，保证了程序包的精简性。

\hypertarget{ux754cux9762ux5e03ux5c40}{%
\subsubsection{5.2 界面布局}\label{ux754cux9762ux5e03ux5c40}}

游戏包含四个主要界面，通过 \texttt{QStackedWidget} 进行页面切换管理：

\begin{enumerate}
\def\labelenumi{\arabic{enumi}.}
\tightlist
\item
  \textbf{主菜单}：大标题与反战副标题构成中心视觉，三个功能按钮（新游戏/继续/退出）垂直排列，底部装饰以反战标语
\item
  \textbf{角色创建}：左侧2×3网格展示六大职业卡片，右侧预览面板显示初始属性、初战方向和兵种特征描述文；选中卡片通过动态属性
  \texttt{selected="true"} + QSS伪类实现红色高亮反馈
\item
  \textbf{游戏主界面}：顶部状态栏同时显示玩家信息（姓名/兵种）、双进度条（生命/士气）和场景标注；中央为大面积文本叙事区；底部左侧生成选项按钮，右侧排列存档/读档/退出控制按钮
\item
  \textbf{存档对话框}：模态对话框展示4个存档槽位卡片，通过
  \texttt{SaveInfo}
  结构体读取摘要信息并格式化显示，选中的槽位以红色边框高亮
\end{enumerate}

\hypertarget{ux9009ux9879ux6309ux94aeux6837ux5f0fux533aux5206}{%
\subsubsection{5.3
选项按钮样式区分}\label{ux9009ux9879ux6309ux94aeux6837ux5f0fux533aux5206}}

在游戏主界面中，不同类型的选项使用不同的按钮样式进行视觉区分： -
普通选项（\texttt{choiceOptionBtn}）：深色边框，左侧文字对齐 -
战斗选项（\texttt{combatOptionBtn}）：红色调背景与边框，红色文字加粗，明确提示风险
- 控制按钮（\texttt{gameCtrlBtn}）：小字体，右侧栏排列

这种通过CSS类名进行语义化样式管理的方式，使得UI代码与样式代码解耦，便于主题调整。

\begin{center}\rule{0.5\linewidth}{0.5pt}\end{center}

\hypertarget{ux6784ux5efaux4e0eux6d4bux8bd5}{%
\subsection{6. 构建与测试}\label{ux6784ux5efaux4e0eux6d4bux8bd5}}

\hypertarget{ux6784ux5efaux914dux7f6e}{%
\subsubsection{6.1 构建配置}\label{ux6784ux5efaux914dux7f6e}}

项目使用CMake作为构建系统，\texttt{CMakeLists.txt}
约80行，主要配置包括：

\begin{itemize}
\tightlist
\item
  设定C++17标准和Qt6 AUTO处理（AUTOMOC/AUTORCC/AUTOUIC）
\item
  查找Qt6的Core、Widgets、Multimedia三个必需组件
\item
  声明28个源文件/头文件并注册Qt资源文件
\item
  配置构建后复制音乐资源目录到可执行文件目录
\item
  设置输出名称为``钢铁意志\_第三帝国''
\end{itemize}

前置依赖为Qt 6（含Multimedia模块）、CMake 3.16+和支持C++17的编译器（MSVC
2019+、GCC 9+或Clang 10+）。

\hypertarget{ux529fux80fdux6d4bux8bd5}{%
\subsubsection{6.2 功能测试}\label{ux529fux80fdux6d4bux8bd5}}

项目通过以下方式进行了功能验证：

\begin{enumerate}
\def\labelenumi{\arabic{enumi}.}
\tightlist
\item
  \textbf{各场景节点完整性测试}：逐一验证每个战役的每个节点文本渲染正确、选项按钮生成正确、职业限制和标志条件过滤正确
\item
  \textbf{骰子系统概率验证}：统计多次投掷结果，确认基础范围、加成效果和钳制行为符合设计
\item
  \textbf{存档/读档循环测试}：在游戏各阶段执行存档后退出并读档，验证玩家状态（姓名、职业、HP、士气、标志集合、所在节点）完整恢复
\item
  \textbf{战役间跳转测试}：验证从黄色方案到柏林战役的完整流程中，战役解锁、节点导航和自动存档逻辑正确
\item
  \textbf{异常处理测试}：验证音乐文件缺失时程序静默运行、存档目录不存在时自动创建、JSON数据残缺时使用默认值容错恢复
\end{enumerate}

\begin{center}\rule{0.5\linewidth}{0.5pt}\end{center}

\hypertarget{ux603bux7ed3ux4e0eux5c55ux671b}{%
\subsection{7. 总结与展望}\label{ux603bux7ed3ux4e0eux5c55ux671b}}

\hypertarget{ux9879ux76eeux6536ux83b7}{%
\subsubsection{7.1 项目收获}\label{ux9879ux76eeux6536ux83b7}}

通过本项目的完整开发实践，笔者在以下方面获得了实质性提升：

\begin{enumerate}
\def\labelenumi{\arabic{enumi}.}
\tightlist
\item
  \textbf{面向对象设计能力}：从抽象基类设计到具体场景继承实现，深入理解了继承、多态和接口隔离原则在实际项目中的应用
\item
  \textbf{Qt框架掌握}：熟练运用信号槽机制实现松耦合通信，掌握QSS样式系统实现工业级UI，理解Qt父子内存管理模型的资源生命周期
\item
  \textbf{软件架构意识}：实践了MVC分层架构，将数据定义、逻辑处理和UI展示分离，使得500行以上的代码库保持清晰的结构
\item
  \textbf{系统工程能力}：从CMake构建配置、资源管理到发布打包，完整体验了一个桌面应用程序从开发到交付的全流程
\end{enumerate}

\hypertarget{ux5c40ux9650ux4e0eux6539ux8fdb}{%
\subsubsection{7.2 局限与改进}\label{ux5c40ux9650ux4e0eux6539ux8fdb}}

当前项目存在以下局限，也对应着未来的改进方向：

\begin{enumerate}
\def\labelenumi{\arabic{enumi}.}
\tightlist
\item
  \textbf{无国际化支持}：所有文本硬编码为中文，可利用Qt
  Linguist（\texttt{.ts}文件）实现多语言翻译
\item
  \textbf{音乐文件依赖}：因版权原因音乐文件未随源码分发，可提供静默模式或使用程序化音频替代
\item
  \textbf{窗口固定尺寸}：当前UI使用固定坐标的透明按钮覆盖层（如角色选择卡片、存档触发按钮），窗口缩放时可能偏移，需重构为自适应布局
\item
  \textbf{成就与统计系统缺失}：玩家的道德选择虽然留下标志，但缺乏统一的后续影响汇总或成就解锁面板
\item
  \textbf{叙事分支深度有限}：每场战役的平均分支深度为3-5层，可增加更多中间态节点使叙事更加细腻
\end{enumerate}

\hypertarget{ux7ed3ux8bed}{%
\subsubsection{7.3 结语}\label{ux7ed3ux8bed}}

《钢铁意志：第三帝国的黄昏》是一个以``反战''为内核的交互式文字冒险游戏，在C++17与Qt6的技术栈上，实现了一个完整的分层架构游戏系统。项目从设计到实现历时数周，总计编写约3000行源代码，涵盖5个历史战役、129个故事节点，以及骰子判定、叙事标志、存档回溯等多项游戏机制。通过这一实践，笔者不仅巩固了面向对象程序设计的理论知识，更获得了从需求分析到软件交付的完整项目经验。

\begin{center}\rule{0.5\linewidth}{0.5pt}\end{center}

\hypertarget{ux53c2ux8003ux6587ux732e}{%
\subsection{参考文献}\label{ux53c2ux8003ux6587ux732e}}

{[}1{]} William L. Shirer. \emph{The Rise and Fall of the Third Reich}.
Simon \& Schuster, 1960.

{[}2{]} Antony Beevor. \emph{Stalingrad: The Fateful Siege: 1942-1943}.
Penguin Books, 1999.

{[}3{]} Qt Group. \emph{Qt 6 Documentation}. https://doc.qt.io/qt-6/,
2024.

{[}4{]} Bjarne Stroustrup. \emph{The C++ Programming Language (4th
Edition)}. Addison-Wesley, 2013.

{[}5{]} Erich Gamma, Richard Helm, Ralph Johnson, John Vlissides.
\emph{Design Patterns: Elements of Reusable Object-Oriented Software}.
Addison-Wesley, 1994.

{[}6{]} CMake Project. \emph{CMake 3.16 Documentation}.
https://cmake.org/documentation/, 2024.

\begin{center}\rule{0.5\linewidth}{0.5pt}\end{center}

\emph{愿世界永无战争。}
